% !TeX program = xelatex
% 编译：xelatex EGO4P2_DEPTH_TO_TRAJECTORY_GUIDE.tex（两次）
% 来源：同目录下的 Markdown 讲解稿。样稿仅用作版式参考。
\documentclass[12pt,a4paper]{article}
\usepackage[UTF8,scheme=chinese,fontset=fandol]{ctex}
\usepackage[margin=24mm]{geometry}
\usepackage{amsmath,amssymb,bm}
\usepackage{array,booktabs,longtable}
\usepackage{enumitem}
\usepackage{xcolor,tikz}
\usetikzlibrary{positioning,arrows.meta}
\usepackage{xurl}
\usepackage[hidelinks]{hyperref}

\definecolor{reviewOrange}{RGB}{255,105,0}
\definecolor{reviewRed}{RGB}{170,0,0}
\definecolor{reviewPurple}{RGB}{128,0,128}
\definecolor{reviewGreen}{RGB}{0,180,0}
\colorlet{replycolor}{black}

\linespread{1.12}
\usepackage{indentfirst}
\setlength{\parindent}{2\ccwd}
\setlength{\parskip}{0.45em plus 0.12em minus 0.1em}
\setlength{\emergencystretch}{2em}
\setlength{\LTleft}{0pt}
\setlength{\LTright}{0pt}
\renewcommand{\arraystretch}{1.25}
\setlist{leftmargin=*,itemsep=0.3em,topsep=0.4em}
\setcounter{tocdepth}{2}

\newcounter{topic}
\numberwithin{equation}{topic}
\newcommand{\topicgroup}[2]{%
  \bigskip\par\noindent{\color{#2}\rule{\linewidth}{0.6pt}}\par
  \section*{\textcolor{#2!70!black}{#1}}%
  \addcontentsline{toc}{section}{#1}%
  \colorlet{replycolor}{#2!70!black}%
}
\newenvironment{point}[1]{%
  \refstepcounter{topic}%
  \subsection*{\thetopic\quad #1}%
  \addcontentsline{toc}{subsection}{\thetopic\space #1}%
}{\par}
\newenvironment{reply}{%
  \par\noindent\textcolor{replycolor}{\sffamily\bfseries 答复}\par\nobreak\ignorespaces
}{\par\medskip}
\newcommand{\sourceref}[2]{\hyperlink{source-#1}{#2\textsuperscript{[#1]}}}
\newcommand{\sourceentry}[2]{\item[\hypertarget{source-#1}{[#1]}] #2}

\begin{document}
\section*{EGO4P2 从深度图到飞行路线}

\tableofcontents
\clearpage
\topicgroup{一、读取相机看到的距离}{reviewOrange}
\begin{point}{无人机根据眼前信息选择路线}
\end{point}
\begin{enumerate}
\item \textbf{深度图观测方法：}四张深度图分别告诉系统视野里首先看到了什么。把信息投到一个以机体为中心的单位球面上，得到可以用于规划的方向区域。
\item \textbf{连续改变运动方向方法：} 在这个区域里给每个方向格分配一个势值，让势值从低处过渡到出口。沿势值上升的方向逐小步前进，得到从当前航向出发的方向曲线 \(Q\)。
\item \textbf{实际飞行下一步速度：} 给 \(Q\) 安排每米可以转多少角度，把方向当成曲线的切向，再累加成三维位置。用原始深度图检查这条曲线附近的连续空间，最后依据实际运动状态计算速度命令候选。
\end{enumerate}

图~\ref{fig:chain} 展示四路观测怎样进入求场、参考生成和控制。
\begin{figure}[!htbp]
\centering
\begin{tikzpicture}[
  box/.style={draw=reviewOrange!48!black,rounded corners=2pt,
              fill=reviewOrange!6,align=center,text width=55mm,
              minimum height=8mm,inner sep=4pt},
  flow/.style={-{Stealth[length=2.2mm]},semithick},
  node distance=5mm
]
\node[box] (depth) {四张距离图\\相机参数与拍摄位置朝向};
\node[box,below=of depth] (prep) {补异常值、整理小块\\合成周围方向信息};
\node[box,below=of prep] (domain) {球面可用方向域 \(\Omega\)};
\node[box,right=9mm of depth] (state) {实际状态与期望方向};
\node[box,right=9mm of domain] (field) {粗算、细算\\得到方向曲线 \(Q\)};
\node[box,below=of field] (spatial) {安排每米转多少角度\\累加出空间位置};
\node[box,below=of spatial] (check) {三维参考 \(\gamma\) 与原图净空检查};
\node[box,below=of check] (cmd) {纠正偏离、更新空间检查\\计算下一步速度候选};
\draw[flow] (depth) -- (prep);
\draw[flow] (prep) -- (domain);
\draw[flow] (domain) -- (field);
\draw[flow] (state) -- (field);
\draw[flow] (field) -- (spatial);
\draw[flow] (spatial) -- (check);
\draw[flow] (check) -- (cmd);
\end{tikzpicture}
\caption{EGO4P2 的局部导航流程}
\label{fig:chain}
\end{figure}

后面会反复出现三个名字，可以先这样记。
\begin{itemize}
\item \textbf{势 \(\phi\)：每个方向格上的一个数。} 这是程序算出来的引导数值。它不表示真实高度，也不表示离障碍有多远。
\item \textbf{方向场 \(g\)：每个方向格上的一支转向箭头。} 箭头由势值的变化算出，告诉程序从这个方向接下来往哪边转。
\item \textbf{空间参考 \(\gamma\)：希望机体中心沿着走的一条空间曲线。} 它给出以米为单位的位置，曲线上的任意进度都可以查询。
\end{itemize}


\subsubsection*{开始前先认识几个词}
\textbf{位置}说明“在哪里”，\textbf{方向}说明“朝哪边”。向北飞只说明方向；从房间哪一处开始、经过哪一处，才是位置。\textbf{速度}还包含快慢。\textbf{速率}只表示快慢，不含方向。

\textbf{坐标系}是一套描述位置和方向的参照。世界坐标系固定在环境中；机体坐标系跟着无人机转动；相机坐标系跟着镜头。相同的“向北”，相对机头可能是向前，也可能是向左，所以不同坐标系的数据要先换到同一套参照中才能比较。

\textbf{向量}可以想成一支箭头，有方向，也有长度。\textbf{单位向量}把长度规定为 1，便于只表达方向。\textbf{归一化}就是保留箭头方向，把它的长度改为 1。零长度箭头没有确定方向，不能这样处理。

\textbf{矩阵}是一张按行、列排列的数字表。它在文中主要用来转换坐标，或把许多方程合在一起计算。\textbf{旋转矩阵}用于换算朝向；右上角的转置符号表示行、列互换。对这里的旋转矩阵，转置可以完成反方向的坐标转换。

\textbf{网格}是把一个区域分成许多小格；\textbf{节点}是参与计算的格或合并后的格；\textbf{域}就是本次允许计算的区域。\textbf{离散}表示只存有限个格或点；\textbf{连续}表示中间任意位置也有定义。

\textbf{标定}是先确定相机怎样成像、安装在哪里以及朝向哪里。\textbf{投影}是计算空间中的点或物体，会出现在图像的什么位置。\textbf{切平面}可以想成轻贴在球面某一点的一张小纸片；沿这张纸片的方向叫切向，垂直纸片的方向叫法向。换成别的曲面时，也按同样方式理解。

\textbf{导数}描述“变化得有多快”。\textbf{积分}把一路上的微小变化累加起来。例如，知道每小段朝哪边走，把这些小段位移累加，就能得到新位置。

\textbf{参考}是希望跟随的路线；\textbf{候选}是算出来、还需要检查的结果；\textbf{证据}是本次观测中足以支持通行判断的数据。算出了路线，不等于已经确认整条路线都能飞。

\textbf{保守处理}表示宁可把不确定的地方留下，也不把它直接当作安全。\textbf{误差界}是计算时保留的误差上限，后面会把它计入检查范围。

\subsubsection*{公式的读法}
字母只是给一个量取的短名字，不必先记住所有字母。每个公式后面都有按计算顺序写的说明。下标通常用于区分相机、节点、时刻或坐标系，并不都表示相乘。

\texttt{min} 表示选最小值，\texttt{max} 表示选最大值。求和符号表示把指定的数全部加起来；积分符号表示沿一段范围持续累加。双竖线表示向量的长度，单竖线用于绝对值时表示不计正负的大小。字母上方的点表示它每秒怎样变化。

\texttt{sin}、\texttt{cos} 用来在角度和方向分量之间换算；\texttt{arcsin} 是反过来根据比例求角度；\texttt{atan2} 根据两个分量确定角度。度与弧度都是角度单位，180 度等于圆周率个弧度。\texttt{ms} 是毫秒，1000 毫秒等于 1 秒。

下面保留完整公式和程序条件。正文先讲用途，公式给出计算办法，随后说明每一步的意思。

\subsubsection*{符号与单位对照}

\begin{longtable}{@{}>{\raggedright\arraybackslash}p{0.21\linewidth}>{\raggedright\arraybackslash}p{0.75\linewidth}@{}}
\caption{符号与单位约定}\label{tab:symbols}\\
\toprule
\textbf{符号} & \textbf{含义} \\
\midrule
\endfirsthead
\multicolumn{2}{l}{\small 符号与单位约定（续）}\\
\toprule
\textbf{符号} & \textbf{含义} \\
\midrule
\endhead
\midrule
\multicolumn{2}{r}{\small 接下页}\\
\endfoot
\bottomrule
\endlastfoot
\(C,B,W\) & 相机坐标系、机体坐标系、世界坐标系。\(B_o\) 专指拍摄这组深度图时的机体坐标系。 \\
\(R_{BC},R_o\) & 分别把相机方向转到机体系、把观测时的机体方向转到世界系。方向只旋转，点的位置还要加平移。 \\
\(\lVert a\rVert,\ a^\top b,\ a\times b\) & 向量的长度、点积、叉积。两个单位向量的点积等于夹角的余弦，叉积给出垂直于二者的旋转轴方向。 \\
\(S^2,\Omega\) & \(S^2\) 是三维空间中的单位球面。 \(\Omega\) 是其中有足够观测证据、允许参与本次规划的方向格集合。 \\
\(\zeta,\ s,\ t\) & 累计转角、沿空间参考走过的距离、时间。单位分别为弧度、米、秒。 \\
\(V,L,r_{\rm env}\) & 正常跟踪目标速度 \(12\,\mathrm{m/s}\)，局部观测范围 \(4\,\mathrm m\)，占用包络半径 \(0.30\,\mathrm m\)。 \\
\end{longtable}

公式用弧度进行三角计算。网格参数和讲解示例也用度表示角度，以便直观理解。 \sourceref{S03}{基础参数} · \sourceref{S04}{导航参数}

\begin{point}{四路深度图处理方法}
\end{point}
\subsubsection*{2.1 深度图、时间戳和相机标定}

\textbf{深度图}给像素记录距离。本项目输入的 \(Z\) 是沿相机光轴的深度，单位是米。相机光轴深度是物体点沿相机正前方方向的距离，通常记作 \(Z\)。可以把光轴想成从相机镜头中心向画面正中间伸出去的一条直线。假设相机正对一面墙，墙距离相机正前方 2 米，那么墙上各点的光轴深度都是 \(Z=2\) 米。同一个 \(Z\) 下，画面边缘的物体沿视线量过去会比画面中心更远。

四台相机的图像必须来自同一个仿真时刻，才能合成一组。图像上记录的拍摄时刻叫\textbf{时间戳}。程序还要等到这一时刻的机体位置和朝向，合起来叫\textbf{位姿}，由记录运动状态的里程计消息提供。

拍照后无人机还可能继续运动。因此，程序另外读取现在的位置和速度，不拿拍照时的速度冒充现在的速度。

每组输入还附带身份信息：组号、相机时间戳、标定版本、配置版本、当前仿真时间、参考版本和失效原因。\textbf{版本}用来区分是不是同一套相机参数、配置或参考。\textbf{接口}就是这些数据交给后续程序时约定好的格式。

下面任一情况出现，程序都会拒绝这组输入：缺少相机图像、相机时间戳不同、缺少拍摄位姿、标定或配置对不上、数据时间比现在还晚，或者数据已经过期。观测和实际状态允许的最大时效均为 10 ms 仿真时间，也就是百分之一秒。

\subsubsection*{2.2 把像素变成射线}

先确定每个像素是沿哪个方向看出去的。\textbf{内参}是一台相机自身的成像参数，例如焦距和画面中心，放在数字表 \(K_c\) 中；\(c\) 是相机编号。\textbf{外参}记录相机装在机体什么位置、朝向哪里。

对图像中的像素位置 \((u,v)\)，程序按下面四式依次计算视线、空间点、视线距离和机体系方向：

\begin{equation}
\label{eq:ray}
b_c=K_c^{-1}[u,v,1]^\top,\qquad
P_C=Zb_c,\qquad
\rho=Z\|b_c\|,\qquad
d_B=\frac{R_{BC}b_c}{\|b_c\|}.
\end{equation}

\noindent\textbf{这组公式的意思}\par\nobreak
\begin{enumerate}
\item 第一式把像素的横、纵位置换成从镜头看出去的一条视线。内参的逆矩阵用来做这种反向换算；方括号中的 1 是这种写法需要的第三个分量。
\item 第二式用光轴深度乘视线系数，得到相机眼中的三维点。
\item 第三式求镜头到该点的直线距离。画面越靠边，这个距离越可能大于光轴深度。
\item 第四式保留视线方向、把长度变为 1，再换算到机体坐标系，供方向球面使用。
\end{enumerate}


\(P_C\) 是相机系中的三维点，\(\rho\) 是沿该射线的距离，\(d_B\) 是长度为 1 的机体系方向。举例来说，如果 \(b_c=(0.5,0,1)\) 且 \(Z=2\,\mathrm m\)，则 \(P_C=(1,0,2)\,\mathrm m\)，而射线距离为 \(\sqrt5\,\mathrm m\)。

为了快速把四张图放到同一个方向球面上，求场时先作一个近似：保留四台相机各自的朝向，但把它们的镜头中心都当作位于机体中心。这样拼出的方向信息叫\textbf{虚拟共心环视}。真实安装位置之间的偏移，在这一步暂不计入。

在球面上，用方位角 \(\alpha\) 表示往水平面哪边转，用仰角 \(\beta\) 表示向上抬或向下俯多少。两个角可以换算成三个方向分量：

\begin{equation}
\label{eq:sphere}
d_B(\alpha,\beta)=
(\cos\beta\cos\alpha,\ \cos\beta\sin\alpha,\ \sin\beta).
\end{equation}

\noindent\textbf{这组公式的意思}\par\nobreak
括号内三个数依次表示方向在三个坐标轴上的分量。方位角决定水平朝哪边，仰角决定上抬或下俯多少；正弦、余弦负责把这两个角换成方向箭头。这支箭头长度为 1，只表示朝向。


相机参数不变时，同一个像素的视线方向也不变。因此，程序提前算好一张对应表：球面上一个格的中心，落在第几台相机的哪个像素块中。运行时查表即可。

这里先用近似模型快速筛方向。后面要检查无人机整段飞行会占据的空间，必须回到\textbf{相机真实的安装位置和朝向}。否则，同一个空间位置可能被错投到别的像素上，查到错误的距离。

\subsubsection*{2.3 原始图、补值图和两级摘要}

有些像素没有给出正常距离。NaN 表示不是一个可用数值；无穷、零和负值在这一步也按异常处理。程序只修补这些像素，正常像素保持原样。

修补时，看\textbf{原始图}中以该像素为中心的 \(3\times3\) 小窗口。只选其中有限、并且大于零的深度值。把这些数加起来，再除以它们的个数，就是填进去的平均深度。

数学上，用 \(\mathcal V_p\) 表示像素 \(p\) 周围这些有效原始像素的集合。只有 \(p\) 异常、并且 \(\mathcal V_p\) 中至少有一个有效像素时，才计算 \(Z_{\rm filled}(p)=\sum_{j\in\mathcal V_p} Z_{\rm raw}(j)/|\mathcal V_p|\)。分子是有效深度的总和，分母是有效像素的个数。

如果周围一个有效值都没有，\(p\) 就继续保持异常。所有修补都只看原图；刚补出的值，不再拿去帮助下一处补值。这样不会把一次猜测继续传开。

四台相机分别交给四个长期保留的工作线程处理。\textbf{线程}是程序中可以单独安排执行的一路工作；\textbf{常驻}表示不用每来一张图就重新建立它。

每路完成四件事：修补异常值、整理像素块信息、准备原尺寸图像的区域查询数据、生成这台相机对应的方向数据。

\textbf{块摘要}就是把一小片像素整理成少量统计数。一级按 \(4\times4\) 像素分块；二级再把相邻的 \(3\times3\) 个一级块合成一块。图像边缘不足整块时，按实际像素处理。各级都保留近障碍需要的最小射线距离，以及原始有效像素数。

二级平均深度要按一级块中补值后有效像素的个数加权。也就是说，包含有效像素多的块，应占更大的份量。空块既不参与求平均，也不参与取最小值：

\begin{equation}
\label{eq:pool}
\bar Z_2=\frac{\sum_{j:\,n_j>0} n_j\bar Z_j}{\sum_j n_j},
\qquad
\rho_{\min,2}=\min_{j:\,n_j>0}\rho_{\min,j}.
\end{equation}

\noindent\textbf{这组公式的意思}\par\nobreak
左式先把每个一级块的平均值乘它的有效像素数，再把这些结果相加，最后除以有效像素总数。这样等价于让每个有效像素贡献相同份量。右式从非空一级块中选最小射线距离，以保留最近障碍的信息。下标 2 表示二级摘要，j 是一级块的编号。


\(n_j\) 是第一个式子使用的计数：第 j 个一级块中，补值后仍为有限正值的像素数。它与\textbf{原始有效像素数}不是同一个统计量。程序另外保存原始有效数、原始最小距离和总像素数。

原图和补值图各有一套区域查询索引。\textbf{索引}是提前整理好的辅助数据，用来快速回答“这个矩形内是否每个原始像素都有效，最浅的深度是多少”。查询仍依据原尺寸图像，没有因为做了块摘要就丢掉原图。

四路完成后，协调线程合并方向信息，再做求场和曲线计算。后面的这些工作不属于四路相机预处理并行。

\begin{point}{从深度图中筛出可用方向}
\end{point}
\subsubsection*{3.1 障碍膨胀和三种状态}
把无人机近似成半径 \(r_{\rm body}=0.10\,\mathrm m\) 的球，再预留 \(0.20\,\mathrm m\) 的跟踪误差预算，得到统一包络 \(r_{\rm env}=0.30\,\mathrm m\)。\textbf{障碍膨胀}可以理解为给障碍画一圈安全边。这样即便机体中心没有撞到物体，机体和预留余量可能触及的方向也会被排除。

从观察中心看，距离 \(\rho>r_{\rm env}\) 的障碍物对应的点周围画一个半径为 \(r_{\rm env}\) 的球，球的视角半径为 \(\arcsin(r_{\rm env}/\rho)\)。代码以 5 cm 为单位分档，在对应档的下界 \(\rho_{\rm lo}\) 上计算：

\begin{equation}
\label{eq:inflation}
\delta_{\rm base}=
\begin{cases}
\pi,&\rho_{\rm lo}\le r_{\rm env},\\
\arcsin(r_{\rm env}/\rho_{\rm lo}),&\rho_{\rm lo}>r_{\rm env}.
\end{cases}
\end{equation}

\noindent\textbf{这组公式的意思}\par\nobreak
这里算的是\textbf{障碍周围要排除多宽的方向范围}，也就是从障碍中心方向向外张开的半角，不是无人机应该转多少度。

“档下界”是这一距离档中最近的距离。例如，距离被分到 2.00 米至 2.05 米这一档，就按 2.00 米查表。同样大小的球越近，看起来越大，排除角度也越大。

分段公式的两行表示两种情况。档下界大于包络半径时，用反正弦计算球的视角半径；档下界不大于包络半径时，代码直接取圆周率，也就是 180 度，作全方向排除。后一行是保守处理规则，不能当作正常远处小球的视角计算。


用每档较近的一端代替真实距离，会把障碍算得稍近一些，因此排除角度会稍大一些。这是有意留余量。

查表时还考虑两种误差。第一，连续方向被放入网格后会有位置偏差，所以加网格量化余量。第二，一个像素块不是一个点，它本身也占一定视角，所以近障碍排除还加上像素块的最大角半径。程序用能包住这些方向的角度窗口做排除标记。只有补值后最小射线距离满足 \(\rho_{\min}\le L=4\,\mathrm m\) 的近块参与这一排除。

多台相机看到同一个方向格时，取所有覆盖块中最小的有效 \(\rho_{\min}\)，保留较近的观测。\textbf{原始完整覆盖}则要求至少有一台相机，对应像素块内的原始深度全部有效。仅仅补值后看起来完整，还不够。随后，每个方向格归入下表中的一种状态。

\begin{longtable}{@{}>{\raggedright\arraybackslash}p{0.21\linewidth}>{\raggedright\arraybackslash}p{0.75\linewidth}@{}}
\caption{球面方向格的三种状态}\label{tab:states}\\
\toprule
\textbf{状态} & \textbf{程序含义} \\
\midrule
\endfirsthead
\multicolumn{2}{l}{\small 球面方向格的三种状态（续）}\\
\toprule
\textbf{状态} & \textbf{程序含义} \\
\midrule
\endhead
\midrule
\multicolumn{2}{r}{\small 接下页}\\
\endfoot
\bottomrule
\endlastfoot
\texttt{Unavailable} & 有深度观测，且该方向被近障碍的膨胀区域排除。 \\
\texttt{Available} & 中心有观测且未被排除。查询窗口内每一格都有原始完整覆盖，融合最小射线距离严格大于 \(L+r_{\rm env}=4.30\,\mathrm m\)。 \\
\texttt{Unknown} & 剩余情况。包括看不到、原始覆盖不完整、深度证据不足等。 \\
\end{longtable}

\(\Omega\) 是通过筛选的方向格名单，只包含 \texttt{Available}。\texttt{Unknown} 表示证据不足，不能当作“没有障碍”。

这一步只是筛选方向。后面把曲线转换为以米表示的位置后，还要用真实相机的原图单独检查。两处使用同一个 \(0.30\,\mathrm m\) 包络定义。空间检查直接看原图，不在已经膨胀过的 \(\Omega\) 上再加一遍包络。 \sourceref{S13}{四路合并和三态判定}

\subsubsection*{3.2 用两个角表示三维方向}

可以把方向球面想成一个地球仪。地球仪上的位置用经度、纬度两个数就能找到；这里用方位角、仰角两个数来找方向。

球面点虽然写成三个坐标，但必须满足 \(x^2+y^2+z^2=1\)，也就是到球心的距离固定为 1。三个坐标不能随意独立改变，所以它是\textbf{二维曲面}。在上面算势，称为二维调和场；向上、向下、向左、向右的三维方向仍然都能表达。

细网格是 \(1^\circ\times1^\circ\)，粗网格以 \(3^\circ\times3^\circ\) 格为基础。它们的单位是角度。上节的 \(4\times4\) 和再 \(3\times3\) 是图像像素块尺度。这两种网格在代码中分别配置。 

\clearpage
\topicgroup{二、从可用方向生成空间路线}{reviewRed}
\begin{point}{给方向格赋值并计算转向指引}
\end{point}
\subsubsection*{4.1 固定出口和边缘的计算规则}

前面已经知道哪些方向可以参与规划。现在要决定从当前方向怎样逐步转到目标方向。程序给每个可用方向格放一个数，叫\textbf{势值}。

这里的\textbf{出口}是本次选定的目标方向，不是门口，也不是一个空间终点。程序把出口附近的格固定为高势 1，把另一侧选定的格固定为低势 0，再计算中间格的数。后面沿势值上升的方向走，就得到一串逐步转向的指引。出口的具体选法见第 4.2 小节。

可以把它想成一张边缘高度已固定的薄膜：中间部分需要平滑地衔接各处的高度。在本项目中，满足局部平衡条件的势称为\textbf{调和势}。在连续模型中，内部势值满足下面的方程：

\begin{equation}
\label{eq:laplace}
\Delta_{S^2}\phi
=\frac{1}{\cos^2\beta}\frac{\partial^2\phi}{\partial\alpha^2}
+\frac{1}{\cos\beta}\frac{\partial}{\partial\beta}
\left(\cos\beta\,\frac{\partial\phi}{\partial\beta}\right)=0.
\end{equation}

\noindent\textbf{这组公式的意思}\par\nobreak
这行公式规定：除固定边界以外，内部每一格的数值变化应与周围保持平衡。它不会要求所有格都等于零。

方位角是左右绕转的角，仰角是上下抬俯的角。第一项描述方位方向上的变化，第二项描述仰角方向上的变化；它们加起来为零。偏导符号表示先只看一个角变化时，势值怎样变化。公式中的余弦用于修正球面不同位置的几何差异。


式中的 \(\Delta_{S^2}\) 称为\textbf{Laplace–Beltrami 算子}。算子就是对一个函数按规定执行的数学运算；这里用它描述球面上每一小块与周围小块之间的平衡。式子等于零，表示内部没有额外的源或汇。

\(\cos\beta\) 用来修正球面不同位置的距离。想象在地球仪上向东转过同样的经度：赤道处走过的距离长，靠近北极时距离短，因为经线在极点汇到一起。

因此，同样是相差一格的方向，在不同仰角处实际靠得并不一样近。公式必须把这种差别算进去。这里说的是\textbf{方向球面上的距离}，不是无人机实际飞了多少米。

\textbf{边界条件}就是规定边缘如何计算。\textbf{Dirichlet 边界}直接固定格子的数，例如出口为 1、低势格为 0。\textbf{零 Neumann 边界}则规定，跨出这条边界的法向通量为零。

\textbf{法向}指垂直边界的方向。\textbf{通量}在这里描述相邻格之间通过边界的数学交换量，不是空气流量。这条规则在代码中通过不向域外格添加连接来实现。

两种规则\textbf{在同一次求场中一起使用}：高、低势格提供固定值；面对障碍、未知区域和求解域外的边缘采用零通量。

程序通过不连接障碍、未知区域和当前求解域之外的格，落实零通量规则。它约束的是势场的计算；后面的方向积分和空间管道检查仍需分别确认曲线没有越界。 \sourceref{S16}{球面权重} · \sourceref{S17}{域外面处理}

\subsubsection*{4.2 选择开始方向和目标方向}

先确定转向从哪里开始。例如，无人机正在向东飞，想转向北方，就应从“向东”开始逐渐改变方向。

这个起点记为 \(q_0\)，是球面上的\textbf{方向}，不是无人机所在的空间位置。\(v_W\) 表示当前世界坐标系中的实际速度，公式为

\begin{equation}
\label{eq:start}
q_0=\frac{R_o^\top v_W}{\|v_W\|}.
\end{equation}

\noindent\textbf{这组公式的意思}\par\nobreak
可以按“换参照、去掉快慢”来读。分子把世界系实际速度转成拍照时机体系中的速度；分母是实际速率。相除后，长度变为 1，只剩当前运动方向。

这个方向标出球面上的积分起点。它不固定该处的势值，也不要求该处就是低势 0。若无人机还在向东运动，就从向东开始转，不能只因指令想向北，就把起点直接放到北方。


当当前速度无法提供有效方向时，程序先尝试上一次有效的世界系速度，再尝试初始期望方向。滚动切换参考时，则优先使用旧参考在连接处的切向，使新参考从连接处的前进方向开始。

公式中使用 \(R_o^\top\)，是为了把世界系方向转回拍摄深度图时的机体坐标系。方向域保存在这个观测坐标系中，起点也必须用同一坐标系表达。

程序先检查 \(q_0\) 所在的格是否可用，再沿方位和仰角的四个相邻格向外搜索。所有能够通过可用格连续走到的格，合起来就是起点所在的\textbf{连通分量}。即使另一个方向看起来没有障碍，如果中间被不可用格隔断，也不能直接把它当作可达出口。

如果期望方向所在格属于这个连通分量，就选它作为出口中心。否则，程序在该分量中选择最接近期望方向的格，并记录这次调整。实现中通过比较单位方向的点积来选择；点积越大，两个方向的夹角越小。

粗场的低势中心选在该分量中离出口方向最远的一侧。高、低势边界各占据一小片球面区域，称为\textbf{边界帽}。它们的角半径取 \(3^\circ\) 与两中心角距的 \(0.24\) 倍中的较小值，标签重叠时高势优先。低势中心与积分起点 \(q_0\) 分别选择，两者并不要求重合。 \sourceref{S18}{起点来源} · \sourceref{S19}{连通域与出口} · \sourceref{S20}{边界帽}

\subsubsection*{4.3 大范围粗算与局部细算}

\begin{itemize}
\item \textbf{粗层先看整个可达范围。}

一个 \(3^\circ\) 粗格原本包含九个 \(1^\circ\) 细格。只有三项条件都满足才合并：二级图像摘要支持它可用；九个细格都属于起点所在的连通分量；九格的边界标签一致。

标签区分普通格、高势出口和低势边界。不能合并的粗格，仍保留其中属于这个连通分量的细节点。两块合并格之间的连接权重，等于穿过它们公共边界的细网格连接权重之和。
\item \textbf{细层只细算起点到出口附近。}

搜索时，程序记住每个格是从哪个格走来的，后者叫父节点。从出口沿这些记录往回找，就能得到一条通往起点的格路径。再向周围可用邻格扩三层，形成局部通道。

默认细场保留粗图的高势出口标签，另选一个远侧格作为低势，通道侧边设为零通量。如果整个通道都被标成出口，就从非起点格中另选低势格。

通道边界是为了节省计算而截取的，并不全是真实障碍。如果积分被这种人工边界挡住，又满足重试条件，就扩大通道再算。细场总共最多尝试三次。
\item \textbf{另有一个可选的边界处理办法。}

程序也支持把对应粗节点的势值，直接用作人工通道边界的势值。这叫粗细耦合。当前默认使用局部零通量分支，没有默认采用这种数值传递；粗场通过可达范围和出口标签影响细场。
\end{itemize}

粗层覆盖较大的方向范围，细层把计算集中在起点通往出口的局部通道。这样可以保留大范围的连通信息，同时减少细网格上需要求解的节点数。这里的\textbf{节点}就是参与计算的格；合并后的粗格也作为一个节点处理。 \sourceref{S21}{粗格聚合} · \sourceref{S22}{局部通道与重试} · \sourceref{S23}{细场边界分支}

\subsubsection*{4.4 用方程算出各格的势值}

计算机不能直接存储球面上无限多个位置的势值，因此要先把连续问题转换成有限个格上的方程。这里采用\textbf{有限体积法}：把每个格看成一个小区域，要求它经过各条边界流出的总通量为零。

相邻格之间的势值差决定通量方向，\textbf{权重}决定这条连接对平衡的影响有多大。方位方向的权重使用格中心仰角 \(\beta_i\)，仰角方向的权重使用两格公共边的仰角 \(\beta_{\rm face}\)。离散平衡方程为

\begin{equation}
\label{eq:flux}
w_\alpha=\frac{\Delta\beta}{\cos\beta_i\,\Delta\alpha},
\qquad
w_\beta=\frac{\cos\beta_{\rm face}\,\Delta\alpha}{\Delta\beta},
\qquad
\sum_{j\in N(i)}w_{ij}(\phi_i-\phi_j)=0.
\end{equation}

\noindent\textbf{这组公式的意思}\par\nobreak
前两式给出两种连接的权重。角增量表示相邻格之间相差多少角度；中心仰角用于方位连接，公共边上的仰角用于仰角连接。

最后一式读作：对格 i 的每个有效邻格 j，算出“自己的势值减去邻格势值”，乘这条连接的权重，再把全部结果相加，要求总和为零。这是内部格的平衡规则。权重越大，这个邻格对平衡的影响越大。


\(N(i)\) 表示与格 \(i\) 相连的域内邻格集合。球面绕一圈会回到原处，因此方位角的 \(-180^\circ/180^\circ\) 两端周期相连；上下极点则不添加跨越极面的邻格。

每个未知格都有这样一条平衡方程。把所有方程放在一起，并把固定边界的已知势值移到等号右边，就得到矩阵形式：

\begin{equation}
\label{eq:matrix}
\begin{gathered}
A\phi=b,\qquad
A_{ii}=\sum_{j\in N(i)}w_{ij},\\
A_{ij}=-w_{ij}\ \text{（相连的未知格）},\qquad
b_i=\sum_{j\in N(i)\cap D}w_{ij}\phi_j,
\end{gathered}
\end{equation}

\noindent\textbf{这组公式的意思}\par\nobreak
\begin{itemize}
\item 左边的大写 A 收集所有连接关系和权重；势值列是要求出的答案；右边的小写 b 收集已知边界带来的数值。
\item 对角项表示一个未知格所有有效连接权重的总和。
\item 两个相连未知格之间填入负的连接权重，表达两格势值相减的关系。
\item 邻格势值若已固定，就不再作为未知量求解，而是把“权重乘固定势值”加到右端。
\end{itemize}
这里的 A 和 b 是方程组中的名字，与前面表示相机视线的小写 b 不是同一个对象。


其中 \(D\) 是已经固定势值的格集合。若同一对节点之间存在多条连接，程序会先把它们的权重相加。大部分节点彼此没有直接连接，对应矩阵元素为零，因此这个矩阵称为\textbf{稀疏矩阵}。

粗层和细层都使用 Eigen 数值库中的稀疏 LDLT 方法求解。\textbf{LDLT 分解}把适用的对称矩阵拆成三角矩阵和对角矩阵，再分步求出未知势值。求解后，程序还检查结果是否为有限数、势值是否落在允许范围内，并核对以下残差条件：

\begin{equation}
\label{eq:residual}
\max_{i\ \text{为未知格}}\frac{|b_i-(A\phi)_i|}{A_{ii}}\le 10^{-7}.
\end{equation}

\noindent\textbf{这组公式的意思}\par\nobreak
对每个未知格，先比较“右端应得到的数”和“把求出的势值代回左边得到的数”，取差的绝对值，再除以该格的对角项。选所有格中最大的一个，要求它不超过千万分之一。检查的是离散方程的计算误差，不是无人机的位置误差。


\textbf{残差}就是把算出的势值代回离散方程后，等号两侧还相差多少。上式先用对应的对角项归一化，再检查所有未知格中的最大差值。

\textbf{缓存}就是把上次结果留着，下次符合条件时直接使用。但不能因为画面“看起来差不多”就复用。

粗层和细层分别核对边界标签、边界数值、求解容差、相关节点连接和权重。\textbf{容差}是允许的计算误差上限。只有严格检查一致，才复用已经验证的解；输入变化则需要重新处理。 \sourceref{S16}{离散权重} · \sourceref{S24}{稀疏装配} · \sourceref{S25}{LDLT 后端} · \sourceref{S26}{严格缓存与残差}

\begin{point}{沿转向指引连成方向曲线}
\end{point}
\subsubsection*{5.1 从格子的数值算出转向箭头}

求场后，每个网格上有了一个势值。接下来需要知道：落在两个网格点之间时，应当朝哪边转。程序先用周围有效格估计角点的势值和导数，再通过二维 \textbf{Hermite 插值}构造连续函数。插值就是根据已有离散数据估计中间位置；Hermite 方法同时利用数值和斜率，使连接更平滑。

在普通区域，程序对插值得到的势函数 \(\tilde\phi\) 求\textbf{球面梯度}。梯度指向势值增长最快的方向，其公式为

\begin{equation}
\label{eq:gradient}
g(d)=\nabla_{S^2}\tilde\phi(d)
=\frac{\partial_\alpha\tilde\phi}{\cos\beta}e_\alpha
+\partial_\beta\tilde\phi\,e_\beta,\qquad
\hat g(d)=\frac{g(d)}{\|g(d)\|}.
\end{equation}

\noindent\textbf{这组公式的意思}\par\nobreak
分别算出势值沿左右、上下变化得有多快，再沿这两个球面切向组合成箭头。方位方向要除以仰角余弦，修正它在球面上的实际长度。最后除以箭头自身的长度，得到单位箭头。帽子符号表示归一化后的结果。

箭头贴着方向球面，告诉程序“运动方向往哪边转”；它不是直接发送给无人机的空间速度。


\(e_\alpha,e_\beta\) 分别是沿方位角和仰角变化的单位切向，可以看作贴在球面上的两个局部方向。梯度 \(g(d)\) 垂直于球面半径 \(d\)，因此它指示的是沿球面往哪里转。出口势值较高，所以积分沿势值上升的一侧进行。只有梯度模长大于 \(10^{-7}\) 时，程序才把它归一化为单位方向。

“调和”与“谐波”都是英文 \texttt{harmonic} 的译法。这里先求标量势，再对势求导得到方向场；粗场和细场则是两种计算层级。前面的残差检查负责核对离散方程，插值负责为后面的积分提供连续查询。

靠近南北极时，\(\cos\beta\) 很小，直接按方位角计算容易不稳定。程序在 \(|\beta|>84^\circ\) 时换用极点附近一整片球面数据，这片区域叫\textbf{极冠}。

它先找一个尽量符合这些数据的简单函数：\(\phi_{\rm cap}=a_0+a^\top(d-d_{\rm pole})\)。这种“用简单函数贴合已有数据”的做法叫\textbf{拟合}；这里使用常数加线性变化，称为\textbf{仿射势}。再从梯度中去掉沿球面半径的分量，得到贴着球面的箭头 \(g_{\rm cap}=(I-dd^\top)a\)。

为了让普通区域与极区平滑衔接，程序在 \(84^\circ\) 到 \(86^\circ\) 之间用权重 \(w=3u^2-2u^3\) 混合两种势。\(u\) 表示仰角绝对值在这两度过渡范围内的相对位置，从 0 变化到 1。由于权重也随位置变化，对混合势求导时必须包括 \(w\) 的导数：

\begin{equation}
\label{eq:polar}
g=(1-w)g_{\rm H}+w\,g_{\rm cap}
  +\partial_\beta w\,(\phi_{\rm cap}-\phi_{\rm H})e_\beta.
\end{equation}

\noindent\textbf{这组公式的意思}\par\nobreak
三个部分要一起看：普通区域的梯度乘普通区域权重，极区的梯度乘极区权重，再加上“混合权重本身在变化”带来的修正。最后一项不能省掉，否则得到的箭头就不是前面那个混合势的真实梯度。

下标 H 表示普通 Hermite 插值结果，cap 表示极冠结果。极区拟合式中的常数项决定势值基准，系数向量决定局部变化；从向量中去掉朝球心或背离球心的部分，就得到贴着球面的梯度。


如果极冠数据不完整，或者拟合没有成功，程序就无法在该方向给出可信的梯度，查询会返回失败。 \sourceref{S27}{Hermite 插值和极区分支}

\subsubsection*{5.2 沿球面小步转动并检查}

得到方向场后，从起点 \(q_0\) 沿箭头走一小步，再查询新位置的箭头，继续向前。这样逐步连成的线叫作\textbf{积分曲线}。这里走的是单位球面，曲线上的每一个点都代表一个运动方向。累计转角作为进度参数，积分方程为

\begin{equation}
\label{eq:ode}
\frac{\mathrm dQ}{\mathrm d\zeta}=\hat g(Q),
\qquad Q(0)=q_0.
\end{equation}

\noindent\textbf{这组公式的意思}\par\nobreak
左边表示方向曲线每增加一点累计转角，方向如何变化；右边就是当前位置查到的单位转向箭头。后半式规定从前面计算的起始方向出发。

这里的进度用弧度表示。它不是飞行时间，也不是已经飞过多少米。


程序使用\textbf{球面中点法}近似求解。对单位方向 \(q\) 和该处的单位切向 \(v\)，\(\operatorname{Exp}_q(hv)=q\cos h+v\sin h\) 表示沿球面大圆前进角度 \(h\)。\textbf{大圆}是过球心的平面与球面相交得到的圆。

一次计算先沿起点梯度走半步，估计中点位置；再查询中点梯度，把它投回起点的切平面并归一化，用这个方向完成整步。这样既利用了中途的方向变化，也让新点留在球面上：

\begin{equation}
\label{eq:midpoint}
\begin{aligned}
q_m&=\operatorname{Exp}_{q_k}\bigl(\tfrac h2\hat g(q_k)\bigr),\\
v_m&=\frac{(I-q_kq_k^\top)\hat g(q_m)}
           {\|(I-q_kq_k^\top)\hat g(q_m)\|},\\
q_{k+1}&=\operatorname{Exp}_{q_k}(h v_m).
\end{aligned}
\end{equation}

\noindent\textbf{这组公式的意思}\par\nobreak
\begin{enumerate}
\item 第一行先沿起点箭头转半步，得到中间方向。
\item 第二行查询中间方向的箭头，去掉不属于起点切平面的分量，再把剩余箭头变成单位长度。单位矩阵减去方向外积，在这里起到“去掉径向分量”的作用。
\item 第三行从原起点沿修正后的箭头转整步，得到下一个方向。
\end{enumerate}
k 是这一步的编号，h 是本次尝试转过的角度。Exp 是沿球面大圆转动的运算名称，在这里不是普通的指数计算。


每次初始尝试的角步长最多为 \(0.8^\circ\)。程序不仅检查步长两端，还检查起点到中点、中点到终点、起点到终点这三条球面短弧经过的格。它还要求势值满足 \(\tilde\phi(q_{k+1})>\tilde\phi(q_k)-10^{-12}\)，并具有足够的角进度。检查失败时，程序会尝试把步长减半后重算。

每个成功短弧的角长度用 \(\operatorname{atan2}(\|q_k\times q_{k+1}\|,q_k^\top q_{k+1})\) 计算，再逐段累加为总转角。只有到达出口标记格，整条曲线才被接收。弱梯度、未知区域、障碍和其他失败都可能使积分提前停止。如果起点已经位于出口格，也可能得到转角为零的有效结果。

例如，无人机正在向东飞，期望改向北方。此时 \(Q\) 可以表示“向东、略偏东北、再逐渐向北”这样一列方向。它描述了方向怎样变化；实际在什么位置开始转、转弯占据多少空间，还要由下一步的空间参考构造确定。 \sourceref{S28}{球面积分与接受条件} · \sourceref{S29}{角长度计算}

\begin{point}{把转向过程变成实际空间路线}
\end{point}
\subsubsection*{6.1 把一个个方向平滑地接起来}

方向积分得到的 \(Q\) 是离散点列。为了在任意转向进度查询方向，程序用三次 Hermite 曲线段连接相邻点，并根据相邻方向估计公共节点处的单位切向 \(t_i\)。

对于端点方向为 \(q_i,q_{i+1}\)、球面角距为 \(h_i\) 的一段，同一条三次曲线可以写成 Bézier 形式，其四个控制点为

\begin{equation}
\label{eq:control-points}
(b_0,b_1,b_2,b_3)=
\bigl(q_i,\ q_i+\tfrac{h_i}{3}t_i,\
q_{i+1}-\tfrac{h_i}{3}t_{i+1},\ q_{i+1}\bigr).
\end{equation}

\noindent\textbf{这组公式的意思}\par\nobreak
四个控制点从左到右排列。首尾两个就是已有方向点；中间两个沿首尾切向伸出或退回一小段，用来控制曲线离开、进入端点的朝向。角距离除以 3 是这种三次曲线写法所用的换算系数。

两个中间控制点通常不在球面上；下一式还会把算出的曲线拉回单位球面。


\textbf{Bézier 控制点}用于确定一段曲线的形状。第一个和最后一个控制点确定端点，中间两个控制点决定曲线以什么切向离开起点、进入终点。段内参数 \(u\in[0,1]\) 从 0 走到 1，就表示从这一段的起点走到终点。

程序先计算三次曲线 \(B(u)\)，再把结果归一化，使方向始终落在单位球面上：

\begin{equation}
\label{eq:normalize}
\begin{gathered}
B(u)=(1-u)^3 b_0+3u(1-u)^2b_1+3u^2(1-u)b_2+u^3b_3,
\\
q(u)=\frac{B(u)}{\|B(u)\|}.
\end{gathered}
\end{equation}

\noindent\textbf{这组公式的意思}\par\nobreak
第一行按段内进度给四个控制点分配不同份量，再相加得到三次曲线。进度为 0 时在起点，为 1 时在终点。第二行除以所得向量的长度，使结果重新落在单位球面。

这里的大写 B 是三次曲线的值，小写 q 是归一化后的方向；它们与表示机体坐标系的字母 B 用途不同。


段内参数走过相同比例，并不一定对应相同转角。为了按真实转角查询，程序数值计算角长度积分 \(\int_0^u\|q'(\xi)\|\,\mathrm d\xi\)，并加上前面各段的角长度。查询时，再根据给定的累计角 \(\zeta\) 反查段内参数 \(u\)。

按累计转角来查询的连续方向曲线记为 \(\widehat Q(\zeta)\)。对非零长度段，有 \(\|\widehat Q'(\zeta)\|=1\)。也就是说，给定进度增加多少转角，曲线上就前进同样多的转角，不会因某段原来的参数变化较快而跳得更远。

计算机通过数值求积近似累加角长度，并保留查询误差界。过近的重复方向点会跳过。连接到相反方向而无法确定转法、曲线出现尖点，或者不满足正则条件时，会拒绝插值。这里的正则条件是保证方向及其变化能够正常计算的要求。

平滑连接后的曲线可能偏离原来的积分短弧，所以原点列通过检查还不够。程序会对\textbf{平滑后的整段方向曲线}重新检查 \(\Omega\) 覆盖。它先利用控制点限定每一段可能经过的方向范围，必要时把曲线继续细分，再检查这些范围所覆盖的格。 \sourceref{S30}{连续方向曲线} · \sourceref{S31}{连续角域检查} · \sourceref{S32}{检查接线}

\subsubsection*{6.2 安排转弯有多急和转弯要走多远}

\textbf{切向}表示沿曲线前进时的方向。\textbf{曲率}表示每走一米，这个方向改变得有多快；曲率越大，转弯越急。

以恒速 \(V\) 沿半径 \(R\) 的圆弧运动时，曲率为 \(1/R\)，垂直于运动方向的法向加速度为 \(V^2/R\)。项目采用 \(80\,\mathrm{m/s^2}\) 作为参考曲线的转向设计值，据此确定最大曲率和最小转弯半径：

\begin{equation}
\label{eq:curvature}
\kappa(s)=\frac{\mathrm d\zeta}{\mathrm ds},\qquad
\kappa_{\max}=\frac{80}{V^2}=\frac59\,\mathrm{m^{-1}},
\qquad
R_{\min}=\frac{1}{\kappa_{\max}}=1.80\,\mathrm m.
\end{equation}

\noindent\textbf{这组公式的意思}\par\nobreak
第一式表示每走一米，累计方向角增加多少，这就是曲率。第二式用参考转向设计加速度除以目标速度的平方，得到曲率上限。第三式取它的倒数，得到与最大曲率对应的圆弧半径。

同样是转 90 度，慢慢转需要更长的路，急转需要更大的曲率。这里限制的是参考路线的弯曲程度，尚未证明它周围没有障碍。


转弯开始和结束时，曲率也需要逐渐变化。按 \(V=12\,\mathrm{m/s}\) 的速度和 \(0.10\,\mathrm s\) 的渐变时间设计，达到最大曲率所需的单侧渐变长度为 \(\ell_r=1.20\,\mathrm m\)。

设 \(\zeta_e\) 是\textbf{平滑后方向曲线的累计转角}，\(a\) 是单侧曲率渐变段长度，\(b\) 是中间恒曲率段长度。程序让曲率从零线性增加，转角足够大时保持一段峰值，再线性降回零。对于正转角，具体安排如下：

\begin{equation}
\label{eq:profile}
\begin{gathered}
a=\min\!\left(\ell_r,\sqrt{\frac{\zeta_e\ell_r}{\kappa_{\max}}}\right),
\qquad
\kappa_p=\kappa_{\max}\frac{a}{\ell_r},
\qquad\\
b=\max(0,\zeta_e/\kappa_p-a),\\[3pt]
\kappa(s)=
\begin{cases}
(\kappa_p/a)s,&0\le s<a,\\
\kappa_p,&a\le s<a+b,\\
(\kappa_p/a)(2a+b-s),&a+b\le s<2a+b,\\
0,&s\ge 2a+b.
\end{cases}
\end{gathered}
\end{equation}

\noindent\textbf{这组公式的意思}\par\nobreak
\begin{enumerate}
\item 先算单侧渐变段长度 a：转角大时用设定的完整渐变长度，转角小时缩短。
\item 再算本次峰值曲率：渐变段短，峰值也相应降低。
\item 余下还需要的转角，由中间长度为 b 的恒曲率段完成；不需要时 b 为零。
\item 分段公式依次描述渐增段、恒定段、渐减段，以及转完后曲率为零的直线段。
\end{enumerate}
分段条件说明当前走过的距离落在哪一段。这个 a 是渐变段长度，与其他章节用同一字母表示的拟合系数或相机系数不同。


把每一小段的“每米转角”沿距离累加，就得到总转角 \(\zeta_e\)。因此，曲率随距离变化的图形面积等于累计转角，转向段的总长度为

\begin{equation}
\label{eq:turn-length}
L_{\rm turn}=2a+b=
\begin{cases}
2\sqrt{\zeta_e\ell_r/\kappa_{\max}},&
\zeta_e<\kappa_{\max}\ell_r,\\
\zeta_e/\kappa_{\max}+\ell_r,&
\zeta_e\ge\kappa_{\max}\ell_r.
\end{cases}
\end{equation}

\noindent\textbf{这组公式的意思}\par\nobreak
总长度就是“前渐变段＋中间段＋后渐变段”。转角较小时，两个渐变段直接接在一起；转角较大时，中间再加恒曲率段。判断大小的门槛是最大曲率乘完整渐变长度。恰好到门槛时，两种写法给出的长度相同。

曲率处处非负时，累计转角是沿途所有转动的总量。它不能用起点和终点两个方向的夹角随意替代。


小转角用两个较短的渐变段即可完成，峰值也低于 \(\kappa_{\max}\)，无需先升到最大曲率。小于 \(10^{-12}\) rad 的转角在代码中按零长度转向段处理。

例如，连续方向曲线 \(\widehat Q\) 的累计转角为 \(90^\circ\) 时，转向段长 \(L_{\rm turn}\approx4.027\,\mathrm m\)，以 \(12\,\mathrm{m/s}\) 匀速走完约需 \(0.336\,\mathrm s\)。这里使用的是沿整条方向曲线累加的转角。如果方向先上抬再回落，中途的转动同样要计入，不能只用起终方向之间的夹角代替。 \sourceref{S33}{曲率轮廓} · \sourceref{S04}{设计参数}

\subsubsection*{6.3 沿前进方向累加出位置}

现在已经知道转向的形状，也安排好了每米转多少角度。\textbf{弧长}就是沿曲线走过的距离。走到弧长 \(s\) 时，程序先计算累计转角，再查询对应方向并旋转到世界坐标系，最后沿这些方向积分得到位置：

\begin{equation}
\label{eq:spatial}
\begin{aligned}
\zeta(s)&=\int_0^s \kappa(u)\,\mathrm du,\\
T_W(s)&=R_o\widehat Q(\zeta(s)),\\
\gamma_W(s)&=p_a+\int_0^s T_W(u)\,\mathrm du,\\
\mathbf K_W(s)&=\frac{\mathrm dT_W}{\mathrm ds}
=R_o\widehat Q'(\zeta(s))\,\kappa(s).
\end{aligned}
\end{equation}

\noindent\textbf{这组公式的意思}\par\nobreak
\begin{enumerate}
\item 累加每小段的曲率，得到走到当前距离时总共转了多少角度。
\item 在连续方向曲线上查询这个转角对应的方向，再转到世界坐标系。这就是路线在此处的前进方向，也叫切向。
\item 从空间锚点出发，把各小段的“方向乘小段长度”累加，得到实际位置。
\item 求前进方向沿距离的变化，得到曲率向量。它同时说明朝哪边弯、弯得多急。
\end{enumerate}
积分中的 u 是累加过程使用的临时变量。它不是本节前面三次曲线的段内进度，也不是后面的执行器输入。


这四行公式依次给出转角进度、世界系切向、世界系位置和曲率向量。\(p_a\) 是\textbf{空间锚点}，也就是这条空间曲线从哪里开始；首次生成通常取当前位置，切换时可以取旧参考上的连接位置。\(\mathbf K_W\) 是\textbf{曲率向量}，既表示转弯有多急，也表示朝哪边弯。

位置、切向和曲率来自同一条连续曲线。为便于快速查询，程序建立间隔为 \(0.05\,\mathrm m\) 的位置表，并额外插入曲率分段的连接点。每个区间用 16 点复合中点法近似积分，也就是在区间内取多个小段中点，用中点处的切向累加位移，同时保存位置误差界。查询并不限于表中的节点，节点之间也能计算。

转向结束后，参考沿末端切向继续接直线，正常参考不会因为到达转向段末端就强制停车。\(4\,\mathrm m\) 限制的是当前观测证据范围，\(L_{\rm turn}\) 描述的是完成转向需要走过的距离。较长的曲线可以先构造出来，但其中哪些部分能够用于执行，还要由后续观测检查决定。

如果只演示从锚点出发、按恒速前进的名义轨迹，可以写成 \(\gamma_W(Vt)\)。实际跟踪时，进度 \(s\) 由当前状态和曲线投影更新，不能只按时间向前累加。 \sourceref{S34}{位置表与连接点} · \sourceref{S35}{位置、直线尾段与曲率查询} · \sourceref{S36}{积分和误差界}

\clearpage
\topicgroup{三、检查通行空间并计算速度}{reviewPurple}
\begin{point}{检查整架无人机能否通过}
\end{point}
\subsubsection*{7.1 用球包住一小段飞行空间}

机体中心能沿一条曲线通过，并不代表整个机体都有足够空间。这里的\textbf{占用管道}是沿参考曲线放置半径为 \(r_{\rm env}\) 的球后形成的空间区域，它同时计入机体尺寸和预留的跟踪误差。

检查时把参考分成若干小段。对弧长区间 \([s_0,s_1]\)，令 \(h=s_1-s_0\)，两端的位置为 \(p_0,p_1\)。程序用一个更大的\textbf{覆盖球}包住这段连续管道，球心和半径计算为

\begin{equation}
\label{eq:ball}
c_{\rm cover}=\frac{p_0+p_1}{2},\qquad
r_{\rm cover}
=r_{\rm env}+\frac{\|p_1-p_0\|}{2}
+\frac{K_{\rm bound}h^2}{8}+\epsilon_{\rm pos}.
\end{equation}

\noindent\textbf{这组公式的意思}\par\nobreak
球心放在两个端点的中间。半径依次加四部分：机体与跟踪误差的包络、端点连线长度的一半、曲线可能偏离连线的最大距离，以及位置计算误差。

这样检查一个球，就能同时包住一整小段路线及其周围管道。只检查两端点是不够的，因为中间可能经过障碍。


\textbf{弦}是连接曲线两端点的直线段。以上覆盖球放在弦的中点，半径中的半弦长用于包住整条端点连线；\(K_{\rm bound}h^2/8\) 再覆盖曲线偏离直线的部分。\(\epsilon_{\rm pos}\) 用来计入位置计算误差，取两端误差界的较大值并加上舍入余量。机体和跟踪误差所需的包络也已包含在半径中。

计算覆盖球时，\(K_{\rm bound}\) 取 \(\kappa_{\max}\) 乘参考旋转的范数上界。这里的\textbf{范数上界}可以理解为数值变换对向量长度影响的最大估计，程序用它避免低估弯曲程度。

只有把累计角长误差也计入后，仍能确认整小段已进入直线尾部，弯曲项才取零。覆盖球会比真实管道大一点，所以这项检查可能拒绝一些实际上够宽的窄通道。 \sourceref{S37}{连续覆盖球}

\subsubsection*{7.2 回到原始图像检查是否有障碍}

覆盖球描述的是实际空间位置，需要与相机真正拍到的位置对应。因此，这一步使用每台相机的真实安装位置和朝向。程序先把球心从世界坐标系转到拍摄时的机体坐标系，再转到相机坐标系：

\begin{equation}
\label{eq:frames}
c_B=R_o^\top(c_{\rm cover}-p_o),\qquad
c_C=R_{BC}^\top(c_B-t_{BC}),
\end{equation}

\noindent\textbf{这组公式的意思}\par\nobreak
第一式先从球心世界位置中减去拍照时的机体位置，再换到拍照时的机体朝向。第二式再减去相机安装位置，并换到相机的朝向。结果就是“这台相机眼中的球心”。

减位置负责换起点，乘旋转矩阵负责换方向。只有两步都做，才能查到覆盖球真正对应的像素。


其中 \(p_o\) 是拍摄时的机体位置，\(t_{BC}\) 是相机光心在机体中的安装位置。位置变换既要旋转，也要扣除这两个平移量。程序还对数值变换的范数留出余量，以保守地处理球的半径。

若机体系下的球半径为 \(r_B\)，而 \(\|c_B\|+r_B>4\,\mathrm m\)，说明这个球至少有一部分超出了有效观测范围，检查直接返回超范围。

把覆盖球投影到图像上，就能确定需要读取哪些像素。程序用一个矩形包住投影区域，这个矩形称为\textbf{ROI}，即感兴趣区域。

设相机系下的球半径为 \(r_C\)，它沿光轴占据的深度区间为 \([c_{C,z}-r_C,c_{C,z}+r_C]\)，近端和远端分别记作 \(Z_{\rm ball,near}\) 与 \(Z_{\rm ball,far}\)。如果 \(c_{C,z}\le r_C\)，球会触及或跨过相机光心所在的深度平面，程序不进行这里的球投影计算。

投影边缘对应恰好与球相切的视线。对图像横轴或纵轴中的一个，令 \(a\) 为相机内参相应一行的前两个系数，\(k_0\) 为该行的主点偏移。为简化表达，记 \(m=a^\top(c_{C,x},c_{C,y})^\top\)、\(D=c_{C,z}^2-r_C^2\)。该轴上包住球投影的两个边界像素坐标为

\begin{equation}
\label{eq:image-bounds}
\mathrm{pixel}_{\pm}=k_0+
\frac{m c_{C,z}\ \pm\ r_C\sqrt{\max(0,m^2+\|a\|^2D)}}{D}.
\end{equation}

\noindent\textbf{这组公式的意思}\par\nobreak
加号、减号分别求出投影在这一图像坐标轴上的两个边缘。横向算一次得到左右边界，纵向再算一次得到上下边界。主点偏移负责把结果放回图像的像素位置。

根号中的最大值操作，是防止数值舍入把本应非负的量算成很小的负数。分母只有在球完整位于镜头前方时才能按这里的规则使用。此处的小写 a、m、大写 D 只是为投影式临时取的名字，不是渐变长度、米制单位或边界格集合。


横轴和纵轴分别计算一次，再向外取整，就得到需要查询的矩形窗口。如果投影超出画面，程序会裁掉画面外的部分，同时保留“投影是否完整”的标记。\textbf{能查到部分像素，不等于看到了整个覆盖球。}

\begin{itemize}
\item \textbf{确认净空必须同时满足以下条件。}

至少一台相机完整看见覆盖球的投影。这台相机的\textbf{近裁面}必须已知、有限并大于零；近裁面是它能提供有效证据的最近光轴深度，球的近端不能比它更近。

矩形窗口里的\textbf{原始}深度必须全部有效，且满足 \(\min_{\rm ROI}Z_{\rm raw}>Z_{\rm ball,far}\)。这句话读作：窗口里最近的原始深度，也要比球的最远端更远。相等也不算通过。这样才有证据说明需要使用的整个球前面没有被这些观测挡住。
\item \textbf{发现疑似障碍时优先拒绝。}

任何相机只要投影成功，即使只看见球的一部分，也要检查补值图。若窗口内最小补值深度满足 \(\min_{\rm ROI}Z_{\rm filled}\le Z_{\rm ball,far}\)，说明有观测比球的远端更近或与它一样近，程序优先返回疑似障碍。

因此，不能只找到一台相机说“可用”就结束，还要检查其他相机有没有给出否决信息。
\item \textbf{保留未知。} 如果没有相机发现疑似障碍，但也没有任何一路满足完整净空条件，结果就是未知。两台相机各看见一部分，不能在这里拼成完整净空证明。
\end{itemize}

完成这些检查后，需要区分三件事：曲线是否已经构造出来，某一小段是否有完整净空证据，以及下一步实际运动是否已经通过检查。前一项成功不能代替后一项。

全管道诊断会在计算出的观测检查范围内逐段记录状态；滚动执行则优先检查当前位置附近、下一步需要使用的参考前缀。\textbf{参考前缀}在这里指从当前检查起点开始连续得到证据的一段，中间不能跳过未知区间。 \sourceref{S38}{真实相机球投影} · \sourceref{S39}{原图净空与否决规则} · \sourceref{S40}{全管道诊断接线}

\begin{point}{根据实际运动修正下一步速度}
\end{point}
\subsubsection*{8.1 找到路线上的对应位置并纠正偏离}

\textbf{跟踪}就是根据实际位置和速度不断修正运动，使无人机贴近参考曲线。这里介绍的正常跟踪分支用于已接近巡航状态的情况。首次接入要求实际速率与 \(12\,\mathrm{m/s}\) 的差不超过 \(0.001\,\mathrm{m/s}\)，位置偏差不超过 \(0.10\,\mathrm m\)，运动方向偏差不超过 \(5^\circ\)。起步加速和异常停车由其他执行流程处理。 \sourceref{S41}{跟踪接入} · \sourceref{S42}{几何接入检查}

外层先找出“我现在对应路线上的哪里”。它不从头搜整条路线，只在上次进度 \(s_{\rm prev}\) 附近查找。

搜索窗口为 \([\max(0,s_{\rm prev}-0.05),\,\min(s_{\rm cert},s_{\rm prev}+0.25)]\)，单位是米。意思是向后最多看 5 厘米、向前最多看 25 厘米；同时不能退到路线起点之前，也不能超过已经有证据的最远位置。这个最远弧长记为 \(s_{\rm cert}\)。

程序先在位置表中寻找最接近实际位置 \(x\) 的线段，再针对连续曲线最多做一次 \textbf{Newton 修正}。这种方法利用当前偏差和局部变化率估计更准确的位置。令 \(r=x-\gamma_W(s)\)；当最近点位于曲线内部时，位置差应当与切向垂直，即 \(r\cdot T=0\)，由此得到修正式

\begin{equation}
\label{eq:projection}
s_{\rm new}=\operatorname{clamp}_{\rm window}\!
\left(s+\frac{r\cdot T}{1-r\cdot\mathbf K}\right).
\end{equation}

\noindent\textbf{这组公式的意思}\par\nobreak
先根据位置差在切向上的分量，估计沿路线还应往前或往后挪多少；分母考虑路线正在弯曲的影响。把这个修正量加到当前进度，再限制在允许搜索的窗口中。

\texttt{clamp} 表示夹在上下限之间，\texttt{window} 指前文给出的搜索窗口。投影的作用是找到“路线上的哪一点最对应我现在的位置”。


式中的限制操作不允许修正结果跑出搜索窗口。程序还会检查分母能否正常使用、位置差是否已经接近垂直于曲线，以及离曲线有多远。

这里的\textbf{横向误差}指垂直于当前三维切向的偏差，既可能包含左右偏移，也可能包含上下偏移。设 \(I\) 是三维单位矩阵，程序先提取这部分误差，再计算纠偏速度和沿参考前进的速度：

\begin{equation}
\label{eq:cross-track}
\begin{gathered}
e_\perp=(I-TT^\top)r,\qquad
w=-k_e e_\perp,\\
v_d=\sqrt{V^2-\|w\|^2}\,T+w,\qquad
k_e=2\,\mathrm{s^{-1}}.
\end{gathered}
\end{equation}

\noindent\textbf{这组公式的意思}\par\nobreak
第一行从总位置差中去掉沿路线前后的部分，只保留偏离路线的部分。再乘负的纠偏增益，让修正速度朝回到路线的一侧。

第二行把速度分成“沿路线前进”和“回到路线”两部分。两部分互相垂直，所以用平方关系分配它们的大小，使合成速率仍是目标值。增益 2 每秒表示相同偏差下采用多强的拉回作用；它不是把误差在一瞬间消掉。


\(w\) 是指向参考曲线的纠偏速度，与单位切向 \(T\) 垂直。程序根据纠偏速度的大小调整切向分量，使两者合成后的目标速率仍为 \(V\)。这相当于在固定总速度中分出一部分用于回到参考，其余部分用于沿参考前进。

横向误差加上位置误差界必须不超过 \(0.20\,\mathrm m\)。外层更新由累计的实际物理时间触发，目标周期为 5 ms。 \sourceref{S43}{外层更新} · \sourceref{S44}{投影与纠偏公式}

\subsubsection*{8.2 限制方向改变的速度}

外层算出了期望速度，内层还要让当前运动方向平滑地转过去。令 \(d=v_d/V\) 为期望运动方向，\(h=v_k/\|v_k\|\) 为当前实际速度方向，\(\dot d\) 表示期望方向随时间的变化率。程序通过对外层跟踪公式直接求导得到这一变化率，这称为\textbf{解析求导}。

当找到的是曲线内部最近点时，进度每秒改变多少用 \(\dot s=(v_k\cdot T)/(1-r\cdot\mathbf K)\) 计算，切向每秒怎样变用 \(\dot T=\mathbf K\dot s\) 计算。前一式先取实际速度沿路线的部分，再修正弯曲造成的影响；后一式用曲率向量乘前进进度的变化率。

如果最近点仍停在参考起点，而无人机的位置还在起点后方，则设 \(\dot s=0\)。这时不把投影进度当作已经向前走了。

下面的公式按顺序计算位置差、横向误差、纠偏速度和期望方向的变化率。变量上方的点表示对时间求导。记 \(V_\parallel=\sqrt{V^2-\|w\|^2}\) 为沿参考切向分配的速率，则

\begin{equation}
\label{eq:derivatives}
\begin{aligned}
\dot r&=v_k-T\dot s,\\
\dot e_\perp&=\dot r-
    (\dot r\cdot T+r\cdot\dot T)T-(r\cdot T)\dot T,\\
\dot w&=-k_e\dot e_\perp,\qquad
\dot V_\parallel=-\frac{w\cdot\dot w}{V_\parallel},\\
\dot d&=(\dot V_\parallel T+V_\parallel\dot T+\dot w)/V.
\end{aligned}
\end{equation}

\noindent\textbf{这组公式的意思}\par\nobreak
这组式子是在计算“期望方向接下来会怎样变化”，以便提前安排转向。
\begin{enumerate}
\item 实际速度减去参考点沿路线移动的速度，得到位置差变化率。
\item 再扣除沿切向的部分，以及切向自己转动带来的影响，得到横向偏差变化率。
\item 由此算出纠偏速度怎样变，并按固定总速率的要求算出切向速率怎样变。
\item 把这些变化合起来，再除以目标速率，得到期望方向的变化率。
\end{enumerate}
公式里的点积用于取同向分量，变量上方的点都表示随时间变化。这里每行都保留，不能仅看最后一行而忽略前面的依赖。


方向变化还要满足执行限值。\(\operatorname{clip}_M(a)\) 表示\textbf{向量限幅}：方向保持不变，只把超过上限的模长缩小到 \(M\)；没有超限时保留原值。

外层先计算期望角速度，内层再根据实际物理步长 \(\Delta t\) 限制角速度的变化量，最后旋转当前运动方向，得到下一步的速度目标：

\begin{equation}
\label{eq:rotation}
\begin{aligned}
\omega^*&=\operatorname{clip}_{10}\bigl(d\times\dot d
                   +20\,h\times d\bigr),\\
\omega_k&=\omega_{k-1}
       +\operatorname{clip}_{100\Delta t}(\omega^*-\omega_{k-1}),\\
v^*_{k+1}&=V\,\operatorname{Rot}(\omega_k\Delta t)\,h.
\end{aligned}
\end{equation}

\noindent\textbf{这组公式的意思}\par\nobreak
第一行合并“提前跟着参考转”的前馈和“纠正当前朝向偏差”的反馈，再限制转动速度。第二行从上一步角速度出发，只允许在这一物理步内改变有限的量。第三行按允许的角速度旋转当前运动方向，再乘目标速率，得到下一步速度目标。

星号表示目标值；下标 k 表示当前步，k−1 表示上一步，k＋1 表示下一步。叉积在这里给出需要绕哪根轴转动的方向。


第一行包含两部分。\(d\times\dot d\) 是\textbf{前馈}，根据期望方向正在怎样变化，提前给出转向趋势；\(20\,h\times d\) 是\textbf{反馈}，根据当前方向与期望方向的偏差继续纠正。

\(\operatorname{Rot}\) 表示绕角速度向量所指的轴旋转，转角为 \(\|\omega_k\|\Delta t\)，实现采用 Rodrigues 旋转公式。角速度上限为 \(10\,\mathrm{rad/s}\)，角加速度上限为 \(100\,\mathrm{rad/s^2}\)。前者限制方向每秒转得多快，后者限制这种转动速度每秒可以改变多少。 \sourceref{S45}{方向变化率和前馈反馈} · \sourceref{S46}{角加速度限幅和旋转}

\subsubsection*{8.3 补偿执行速度的响应滞后}

执行模型收到速度输入后，实际速度不会立刻变成输入值，而是按响应时间逐步靠近。对于正常非零指令，并且执行加速度限制尚未触发的分支，下一步速度向输入 \(u_k\) 靠近的关系为

\begin{equation}
\label{eq:actuator}
v_{k+1}=v_k+\lambda(u_k-v_k),\qquad
\lambda=\min(1,\Delta t/\tau),\qquad \tau=0.10\,\mathrm s.
\end{equation}

\noindent\textbf{这组公式的意思}\par\nobreak
下一步速度等于“现在的速度＋向输入靠近的一部分”。靠近多少由响应系数决定：实际物理步长除以 0.10 秒，再限制为不超过 1。例如，物理步长为 5 ms 时，未触发其他限制的这一分支每步靠近输入差值的百分之五。

响应时间是执行模型的参数。这条公式只描述前文限定的分支，不能省略“正常非零指令且加速度限制未触发”的条件。


如果直接把速度目标作为输入，模型的响应滞后会使下一步达不到目标。\textbf{逆补偿}利用已知模型反推需要发送的输入：希望下一步达到 \(v^*_{k+1}\)，就把上式反解为

\begin{equation}
\label{eq:compensation}
u_k=v_k+\frac{v^*_{k+1}-v_k}{\lambda}.
\end{equation}

\noindent\textbf{这组公式的意思}\par\nobreak
先算“下一步目标速度减去当前速度”，再除以响应系数，最后加回当前速度。这样得到的内部输入，才有可能让这个执行模型在下一步达到目标。

内部输入与实际速度是两个量。输入可以暂时比目标大，但仍必须通过后面列出的模长、变化率和执行能力检查。


因此，当实际速度目标为 \(12\,\mathrm{m/s}\) 时，内部补偿输入的模长可能超过 12 m/s。这是为了补偿模型响应，并不表示把正常飞行速率目标提高了。

补偿结果仍需通过能力检查。下一步所需的速度变化不能超过 \(150\,\mathrm{m/s^2}\) 乘实际步长；内部输入模长不能超过 \(17\,\mathrm{m/s}\)；相邻输入的变化率不能超过 \(450\,\mathrm{m/s^2}\)。任何一项超限，或者补偿输入接近零，程序都会拒绝这个正常跟踪步。

这里有几组用途不同的加速度参数。\(80\,\mathrm{m/s^2}\) 用于设计参考曲率，\(150\,\mathrm{m/s^2}\) 用于校验执行模型的能力；旧的 \(20\,\mathrm{m/s^2}\) 参考指令限幅则属于另一配置项，不能把它们当成同一个参数。 \sourceref{S46}{逆补偿与拒绝条件} · \sourceref{S47}{模型反馈夹具} · \sourceref{S04}{参数定义}

\subsubsection*{8.4 更新检查结果并决定是否换路线}

\textbf{滚动续证}是随着无人机前进，用新深度图重新检查旧参考前方的一小段。只要原参考仍然适用，就可以继续跟踪，减少重新求场和构造曲线的次数。

每收到一个新深度组，程序先更新图像数据，再从当前进度 \(s\) 向后退 \(0.05\,\mathrm m\)，对齐到位置网格，并保证检查起点不小于零。向前最多检查到 \(s+0.35\,\mathrm m\)；一旦遇到没有净空证据的小段，就在那里截断，不能跳过去继续累积可用长度。

是否继续使用旧路线，要同时看两个条件：期望方向相对这条参考采纳的方向变化不超过 \(1^\circ\)；已经检查通过的连续路线至少覆盖到 \(s+V\Delta t\)，也就是本步按目标速率前进后的位置。

两项都满足时，本周期可以不重新求场、不重新造空间曲线。证据最长只使用 10 ms 仿真时间。若收到的状态时间戳与上次相同，可以更新证据，但不会把同一个控制步执行两次。

滚动层把相邻实际状态的时间戳差作为 \(\Delta t\)，首步采用配置的 5 ms。若间隔超过一个配置物理步、状态过期或时钟回退，程序会进入对应的保护或重置流程。

需要更换参考时，程序先核对连接处的\textbf{位置、切向和曲率向量}。这些检查分别关注是否接在同一位置、前进方向是否衔接、转弯状态是否连续，随后还要核对接入条件和新参考的净空证据。

切换后，程序继续使用累计进度、上一步角速度和补偿输入，不把执行状态全部清零。

新曲线从零曲率开始。如果旧路线此处还在转弯，就不能直接认为二者接得上，需要等到满足连接条件的位置。若旧路线也没有足够的新证据，就进入保护，不能因为新路线暂时接不上而忽略旧路线已经缺证。

参考前方已有净空证据之后，还要检查实际速度候选会把机体带到哪里。程序用 \(p_1=x+v^*_{k+1}\Delta t\) 预测下一步机体中心。\textbf{扫掠空间}是机体沿这一小步运动时占据的全部空间。

程序要求至少一个已经通过净空检查的覆盖球，同时包住本步起点和终点的机体球：

\begin{equation}
\label{eq:sweep}
\max\bigl(\|x-c_{\rm cover}\|,\ \|p_1-c_{\rm cover}\|\bigr)
+r_{\rm body}\le r_{\rm cover}.
\end{equation}

\noindent\textbf{这组公式的意思}\par\nobreak
先分别算当前中心、下一步中心离覆盖球球心多远，取较大的一个，再加机体半径。若总数仍不超过覆盖球半径，两个端点的整个机体都在同一个已检查球内。

这个覆盖结论对应程序预测的两端点连线运动。它并不单凭两个端点，就保证任意绕弯的实际运动都被覆盖。


球是\textbf{凸集}，即球内任意两点之间的线段也在球内。因此，同一个覆盖球满足上述条件时，也能包住机体沿两端点连线运动的整步扫掠。检查失败则返回保护状态。

在这里描述的 ROS \texttt{--track-rolling} 诊断入口中，结果用于记录速度候选和保护状态，运动发布数为零。后面的实验部分会另行说明 R05 专用仿真入口的验证范围。 \sourceref{S48}{续证窗口和切换验收} · \sourceref{S49}{扫掠计算} · \sourceref{S50}{调度与保护} · \sourceref{S51}{ROS 输出}

\clearpage
\topicgroup{四、与论文对照并读懂实验结果}{reviewGreen}
\begin{point}{EGO4P2 与三维 FLARE 的实现对照}
\end{point}
两种方法都先构造调和势，再沿梯度积分。主要差别在于求场的位置：FLARE 在三维位置空间中计算，EGO4P2 在单位方向球面上计算。因此，前者的积分曲线直接表示空间路径，后者还要把方向变化转换成以米为单位的位置曲线。

FLARE 论文第 III 节先膨胀局部障碍，再从三维工作空间中移除这些障碍区域，形成粗、细自由域。该流程没有把 ESDF 列为建域步骤。\textbf{ESDF} 是一种距离表示，给空间位置记录到最近障碍的带符号欧氏距离；论文在讨论其他方法时提及它。 \sourceref{S52}{FLARE 论文第 III 节、图 2–3 与式 (2)–(11)}

\begin{longtable}{@{}>{\raggedright\arraybackslash}p{0.19\linewidth}>{\raggedright\arraybackslash}p{0.37\linewidth}>{\raggedright\arraybackslash}p{0.37\linewidth}@{}}
\caption{三维 FLARE 与 EGO4P2 的实现对照}\label{tab:flare}\\
\toprule
\textbf{对比环节} & \textbf{三维 FLARE} & \textbf{当前 EGO4P2} \\
\midrule
\endfirsthead
\multicolumn{3}{l}{\small 三维 FLARE 与 EGO4P2 的实现对照（续）}\\
\toprule
\textbf{对比环节} & \textbf{三维 FLARE} & \textbf{当前 EGO4P2} \\
\midrule
\endhead
\midrule
\multicolumn{3}{r}{\small 接下页}\\
\endfoot
\bottomrule
\endlastfoot
场所在的位置 & 三维自由空间，变量 \(\xi\) 是一个米制位置。 & 二维单位球面，变量 \(d\) 是一个三维单位方向。 \\
粗细边界 & 粗场外边界编码期望方向，细场自由耦合边界继承粗势值。 & 粗层用低势 0 与出口 1。默认细场用局部通道、独立低势和继承的出口标签。 \\
积分起点 & 空间锚点 \(a_r\)。 & 实际速度方向或连接切向 \(q_0\)。 \\
梯度积分输出 & 直接得到米制空间参考的几何路径。 & 得到累计转角参数下的方向曲线，还要通过曲率安排和位置积分生成空间参考。 \\
跟踪纠偏 & 用穿过锚点的等势面建立横向坐标。 & 在上次进度附近投影到空间曲线，做横向反馈、方向限幅和执行模型补偿。 \\
\end{longtable}

\subsubsection*{9.1 FLARE 在实际空间中求场}

用 \(W^c\) 表示粗工作空间，用 \(\Gamma_{cf}^{\rm free}\) 表示细场与粗场之间可以传递势值的自由耦合边界，用 \(q_{\rm cmd}\) 表示命令方向。下面两行分别给出粗场和细场的方程；每行左侧规定内部平衡，右侧规定边界取值：

\begin{equation}
\label{eq:flare-potential}
\begin{aligned}
\Delta_{\mathbb R^3}\phi^c&=0
&&\text{在粗自由域内},&
\phi^c|_{\partial W^c}&=q_{\rm cmd}^{\top}\xi,\\
\Delta_{\mathbb R^3}\phi^f&=0
&&\text{在细自由域内},&
\phi^f|_{\Gamma_{cf}^{\rm free}}&=\phi^c|_{\Gamma_{cf}^{\rm free}},
\end{aligned}
\end{equation}

\noindent\textbf{这组公式的意思}\par\nobreak
上标 c 表示粗场，f 表示细场。两行左边都要求各自自由区域内部满足调和平衡。粗场外边界的势值，由命令方向与空间位置的点积给出；细场在与粗场相接的自由边界上，直接使用粗场算出的势值。

竖线和其下方的边界名字表示“在这条边界上取值”。这一组是 FLARE 论文的方程，用于对照；EGO4P2 默认细场的边界处理已经在第 4 节单独说明。


粗、细两层面对障碍的边界均满足 \(\nabla\phi^{c,f}\cdot n=0\)。其中 \(n\) 是垂直于边界的法向，梯度在这个方向上的分量为零，表示不穿过障碍边界。

每次更新势场后，论文在该更新区间使用梯度场 \(u_r(\xi)=\nabla_\xi\phi_r^f(\xi)\)，并由空间锚点 \(a_r\) 选定穿过该点的一条积分曲线。在梯度非零处，把梯度归一化只改变沿线推进的参数速度，不改变几何路径。EGO4P2 在方向球面上采用相似的积分过程，但还需经过第 6 节 的空间积分，才能得到实际位置。 \sourceref{S52}{FLARE 式 (2)–(6)}

\subsubsection*{9.2 FLARE 沿流线描述偏离并纠正}

FLARE 也需要知道无人机偏离参考有多少。它用\textbf{等势面}作定位依据，也就是把所有势值相同的空间点连成一个面。

想象几条相邻路线各自在这个面上留下一个交点。比较交点，就能区分无人机跟着哪条路线走，以及相对目标路线偏到哪边。这种沿一束流线建立局部坐标的办法，叫作 \textbf{flow-box}。它只在相应局部条件成立时使用。

具体做法是沿当前位置所在的流线，追溯到穿过锚点的等势面，再用交点在该面上的两个坐标描述横向位置。论文用 \(\Psi_r(\xi)\) 表示从空间点 \(\xi\) 所在流线找到的交点。\(E_r\) 的两列是该等势面在锚点处切平面的两个正交单位方向，用它们把三维位置差转换成两个横向坐标：

\begin{equation}
\label{eq:flow-box}
\begin{gathered}
\chi_r(\xi)=E_r^\top\bigl(\Psi_r(\xi)-a_r\bigr),\qquad\\
\eta_r(t)=\chi_r(x(t)),\qquad J_r=D\chi_r.
\end{gathered}
\end{equation}

\noindent\textbf{这组公式的意思}\par\nobreak
第一式把“流线与等势面的交点减去锚点”投到两个横向基准方向上，得到两个横向坐标。第二行把无人机当前位置代进去，得到当前横向误差；再计算位置变化会怎样改变这两个坐标。

基准方向的两支单位箭头互相垂直，相当于在这个局部平面上画了两根尺子。这里的大写 D 表示求导运算，与前面的投影分母不是同一个量。


\(J_r\) 是\textbf{Jacobian 矩阵}，中文常称雅可比矩阵，记录空间位置发生微小变化时，两个横向坐标分别会改变多少。\(J_r^\dagger\) 是这里使用的右伪逆，用于把需要的横向修正转换回三维速度。

令 \(K_f\) 为论文中的正定反馈增益矩阵，用来设定各横向误差的纠正强度，并令 \(\hat u_r=u_r/\|u_r\|\) 表示单位梯度方向，速度指令为

\begin{equation}
\label{eq:flare-feedback}
v_{\rm cmd}=v_t\hat u_r-J_r^\dagger K_f\eta_r,\qquad
J_r^\dagger=J_r^\top(J_rJ_r^\top)^{-1}.
\end{equation}

\noindent\textbf{这组公式的意思}\par\nobreak
前半项让无人机沿单位梯度方向前进，后半项根据横向误差把速度拉回参考。反馈矩阵决定拉回的强度；右伪逆把两个横向修正量换成一个三维速度修正。

右边给出这里使用的右伪逆写法，需要相应矩阵可逆，也就是两项横向坐标在局部足够独立。这里的前进速度由论文自己的符号表示，不能直接当成当前项目的固定 12 m/s。


当局部坐标有效、参考保持固定，并采用理想速度执行模型时，论文得到 \(\dot\eta_r=-K_f\eta_r\)。它表示横向误差按设定的反馈规律衰减，这个结论依赖上述条件。

EGO4P2 使用第 8 节 中定义的投影误差 \(e_\perp\)，并考虑离散执行、方向变化限值和响应补偿。两者都安排沿参考前进和横向纠偏，但误差的定义与控制方程不同，不能直接把论文的误差结论作为当前实现的验证结果。 \sourceref{S52}{FLARE 式 (7)–(11)}

\begin{point}{已有实验结果与验证范围}
\end{point}
阅读实验结果时，需要先区分测试对象和计时范围。\textbf{闭环}表示执行后的实际位置、速度会返回控制器，参与下一步命令的计算。模型闭环可以检查控制逻辑与指定模型是否配合，仿真飞行则还包含传感器、通信和场景交互。

\textbf{FIFO5} 表示相关测试中的五个算法线程使用 Linux 的 \texttt{SCHED\_FIFO} 调度策略、优先级为 5，四路相机线程与协调线程分别绑定 CPU。它描述线程怎样获得处理器时间，本身不代表每次计算都能满足时限。

\textbf{P99} 是 99\% 样本位于其下或与之相等的耗时统计位置，用于观察较慢的那部分运行。\textbf{最大值}则是这批样本中最慢的一次。汇报这些指标时，还要说明样本来源、计时包含的阶段，以及超过 5 ms 预算的次数。

\begin{longtable}{@{}>{\raggedright\arraybackslash}p{0.21\linewidth}>{\raggedright\arraybackslash}p{0.40\linewidth}>{\raggedright\arraybackslash}p{0.32\linewidth}@{}}
\caption{模型、ROS 回放和 Isaac Sim 的验证范围}\label{tab:experiments}\\
\toprule
\textbf{证据} & \textbf{已有结果} & \textbf{这组证据覆盖的范围} \\
\midrule
\endfirsthead
\multicolumn{3}{l}{\small 模型、ROS 回放和 Isaac Sim 的验证范围（续）}\\
\toprule
\textbf{证据} & \textbf{已有结果} & \textbf{这组证据覆盖的范围} \\
\midrule
\endhead
\midrule
\multicolumn{3}{r}{\small 接下页}\\
\endfoot
\bottomrule
\endlastfoot
第 5 步模型闭环与 FIFO5 & 8,400 个模型周期中，8,328 个跳过求场，续证加跟踪平均 \(0.858\) ms，P99 \(1.230\) ms，最大 \(1.506\) ms。72 个含规划周期平均 \(4.293\) ms，最大 \(6.467\) ms，12 次超过 5 ms。 & 模型用几何观测谓词提供净空证据。12 次超时均在首次规划。 \\
ROS 原始深度回放 & 实测接收 10 个完整组，经历一次时钟回退。近场原图净空不足时进入保护，运动发布数为零。 & 验证输入链、保护和重置行为。 \\
R05 Isaac Sim 隔离引导飞行 & \(41.50\) 仿真秒，\(308.925\,\mathrm m\)，独立检查 \(8{,}300\) 步扫掠零重叠，终点误差 \(0.169\,\mathrm m\)，最终速度为零。 & 使用历史路线做几何引导、四路实时深度、独立碰撞复核，以及实验专用起停和引导控制。只给终点的自主对照未完成。 \\
R05 滚动分支计时 & 552 个有完整分段计时的周期平均 \(1.175\) ms，P99 \(3.618\) ms，最大 \(4.914\) ms，0 次超过 5 ms。 & 五线程普通调度。其余 7,648 个引导控制周期和外层安全复核缺对应的完整分段计时。整轮飞行的 5 ms 算法时限仍待完整测量。 \\
\end{longtable}

R05 中的\textbf{5 ms 是仿真物理步长}，表示仿真状态每次向前推进的时间。\textbf{墙钟时间}是现实中实际经过的时间；该轮平均每个飞行步约需 18.60 ms，包含渲染、通信、控制和独立碰撞复核。

该实验的 Isaac 深度入口还对“射线没有命中物体”作了专门适配。程序先确认各相机的远裁面大于 \(4.5\,\mathrm m\)，才把该渲染器返回的 \texttt{+Inf} 无命中值映射成 \(4.5\,\mathrm m\) 的光轴深度，局部净空查询仍限制在 \(4\,\mathrm m\) 内。这属于实验入口的特定处理；第 2 节 介绍的通用深度预处理仍按其原有规则判断异常像素。

展示视频使用已验收的 R05 飞行轨迹，在原 Isaac Sim 场景中重新渲染。表中的飞行成绩来自原始运行日志和独立审计，视频回放本身不构成一次新的闭环飞行验证。 \sourceref{S53}{第 5 步实现和计时} · \sourceref{S54}{R05 实验、计时和视频}
\clearpage
\section*{源码、论文与实验索引}
\addcontentsline{toc}{section}{源码、论文与实验索引}
以下项目相对路径以 \path|/home/starry/isaac-data/EGO4P2/docs/| 为基准，行号对应原讲解稿核对的源码版本。本文移动到下载目录或 Overleaf 后，路径仍用于定位原项目文件；本次文字整理未重新验证源码版本或实验结果。
\begingroup\small
\begin{description}[style=multiline,leftmargin=3.1em,labelwidth=2.5em,itemsep=0.3em]
\sourceentry{S01}{构建接线}\par \path|../src/new_controller/CMakeLists.txt|，第 21 行
\sourceentry{S02}{当前项目检查点}\par \path|CURRENT_CHECKPOINT.json|
\sourceentry{S03}{基础参数}\par \path|../candidates/ldlt_parallel/new_controller/include/new_controller/base_config.hpp|，第 5 行
\sourceentry{S04}{导航参数}\par \path|../candidates/ldlt_parallel/new_controller/include/new_controller/navigation_config.hpp|，第 5 行
\sourceentry{S05}{四路接收与对时}\par \path|../src/new_controller/src/navigation_input.cpp|，第 24 行
\sourceentry{S06}{输入身份与校验}\par \path|../candidates/ldlt_parallel/new_controller/include/new_controller/navigation_pipeline.hpp|，第 69 行
\sourceentry{S07}{射线查表}\par \path|../candidates/ldlt_parallel/new_controller/include/new_controller/depth_preprocess.hpp|，第 49 行
\sourceentry{S08}{虚拟环视查表}\par \path|../candidates/ldlt_parallel/new_controller/include/new_controller/depth_virtual_panorama.hpp|，第 45 行
\sourceentry{S09}{邻域补值}\par \path|../candidates/ldlt_parallel/new_controller/include/new_controller/depth_preprocess.hpp|，第 19 行
\sourceentry{S10}{两级摘要和区域查询}\par \path|../candidates/ldlt_parallel/new_controller/include/new_controller/depth_observations.hpp|，第 204 行
\sourceentry{S11}{近块判定}\par \path|../candidates/ldlt_parallel/new_controller/include/new_controller/depth_local_field.hpp|，第 20 行
\sourceentry{S12}{分档与角窗口}\par \path|../candidates/ldlt_parallel/new_controller/include/new_controller/depth_virtual_panorama.hpp|，第 60 行
\sourceentry{S13}{四路合并和三态判定}\par \path|../candidates/ldlt_parallel/new_controller/include/new_controller/depth_local_field.hpp|，第 80 行
\sourceentry{S14}{球面网格}\par \path|../candidates/ldlt_parallel/new_controller/include/new_controller/harmonic_solver.hpp|，第 8 行
\sourceentry{S15}{粗格观测来源}\par \path|../candidates/ldlt_parallel/new_controller/include/new_controller/harmonic_planner.hpp|，第 134 行
\sourceentry{S16}{球面权重}\par \path|../candidates/ldlt_parallel/new_controller/include/new_controller/harmonic_solver.hpp|，第 14 行
\sourceentry{S17}{域外面处理}\par \path|../candidates/ldlt_parallel/new_controller/include/new_controller/grounded_system.hpp|，第 19 行
\sourceentry{S18}{起点来源}\par \path|../candidates/ldlt_parallel/new_controller/include/new_controller/navigation_pipeline.hpp|，第 102 行
\sourceentry{S19}{连通域与出口}\par \path|../candidates/ldlt_parallel/new_controller/include/new_controller/harmonic_planner.hpp|，第 51 行
\sourceentry{S20}{边界帽}\par \path|../candidates/ldlt_parallel/new_controller/include/new_controller/harmonic_adaptive.hpp|，第 32 行
\sourceentry{S21}{粗格聚合}\par \path|../candidates/ldlt_parallel/new_controller/include/new_controller/harmonic_adaptive.hpp|，第 41 行
\sourceentry{S22}{局部通道与重试}\par \path|../candidates/ldlt_parallel/new_controller/include/new_controller/harmonic_planner.hpp|，第 70 行
\sourceentry{S23}{细场边界分支}\par \path|../candidates/ldlt_parallel/new_controller/include/new_controller/harmonic_planner.hpp|，第 162 行
\sourceentry{S24}{稀疏装配}\par \path|../candidates/ldlt_parallel/new_controller/include/new_controller/grounded_system.hpp|，第 14 行
\sourceentry{S25}{LDLT 后端}\par \path|../candidates/ldlt_parallel/new_controller/src/sparse_trial_backend.cpp|
\sourceentry{S26}{严格缓存与残差}\par \path|../candidates/ldlt_parallel/new_controller/include/new_controller/trial_field_solver.hpp|，第 78 行
\sourceentry{S27}{Hermite 插值和极区分支}\par \path|../candidates/ldlt_parallel/new_controller/include/new_controller/harmonic_direction_field.hpp|，第 10 行
\sourceentry{S28}{球面积分与接受条件}\par \path|../candidates/ldlt_parallel/new_controller/include/new_controller/harmonic_curve.hpp|，第 84 行
\sourceentry{S29}{角长度计算}\par \path|../candidates/ldlt_parallel/new_controller/include/new_controller/harmonic_curve_evidence.hpp|，第 75 行
\sourceentry{S30}{连续方向曲线}\par \path|../candidates/ldlt_parallel/new_controller/include/new_controller/spatial_reference.hpp|，第 23 行
\sourceentry{S31}{连续角域检查}\par \path|../candidates/ldlt_parallel/new_controller/include/new_controller/reference_occupancy.hpp|，第 20 行
\sourceentry{S32}{检查接线}\par \path|../candidates/ldlt_parallel/new_controller/include/new_controller/navigation_pipeline.hpp|，第 172 行
\sourceentry{S33}{曲率轮廓}\par \path|../candidates/ldlt_parallel/new_controller/include/new_controller/spatial_reference.hpp|，第 161 行
\sourceentry{S34}{位置表与连接点}\par \path|../candidates/ldlt_parallel/new_controller/include/new_controller/spatial_reference.hpp|，第 183 行
\sourceentry{S35}{位置、直线尾段与曲率查询}\par \path|../candidates/ldlt_parallel/new_controller/include/new_controller/spatial_reference.hpp|，第 228 行
\sourceentry{S36}{积分和误差界}\par \path|../candidates/ldlt_parallel/new_controller/include/new_controller/spatial_reference.hpp|，第 296 行
\sourceentry{S37}{连续覆盖球}\par \path|../candidates/ldlt_parallel/new_controller/include/new_controller/reference_occupancy.hpp|，第 87 行
\sourceentry{S38}{真实相机球投影}\par \path|../candidates/ldlt_parallel/new_controller/include/new_controller/depth_local_geometry.hpp|，第 98 行
\sourceentry{S39}{原图净空与否决规则}\par \path|../candidates/ldlt_parallel/new_controller/include/new_controller/reference_occupancy.hpp|，第 58 行
\sourceentry{S40}{全管道诊断接线}\par \path|../candidates/ldlt_parallel/new_controller/include/new_controller/navigation_pipeline.hpp|，第 184 行
\sourceentry{S41}{跟踪接入}\par \path|../candidates/ldlt_parallel/new_controller/include/new_controller/fixed_reference_tracker.hpp|，第 34 行
\sourceentry{S42}{几何接入检查}\par \path|../candidates/ldlt_parallel/new_controller/include/new_controller/spatial_reference.hpp|，第 280 行
\sourceentry{S43}{外层更新}\par \path|../candidates/ldlt_parallel/new_controller/include/new_controller/fixed_reference_tracker.hpp|，第 66 行
\sourceentry{S44}{投影与纠偏公式}\par \path|../candidates/ldlt_parallel/new_controller/include/new_controller/fixed_reference_tracker.hpp|，第 146 行
\sourceentry{S45}{方向变化率和前馈反馈}\par \path|../candidates/ldlt_parallel/new_controller/include/new_controller/fixed_reference_tracker.hpp|，第 182 行
\sourceentry{S46}{角加速度限幅和旋转}\par \path|../candidates/ldlt_parallel/new_controller/include/new_controller/fixed_reference_tracker.hpp|，第 82 行
\sourceentry{S47}{模型反馈夹具}\par \path|../candidates/ldlt_parallel/new_controller/test/fixed_tracking_cases.hpp|，第 44 行
\sourceentry{S48}{续证窗口和切换验收}\par \path|../candidates/ldlt_parallel/new_controller/include/new_controller/rolling_navigation.hpp|，第 29 行
\sourceentry{S49}{扫掠计算}\par \path|../candidates/ldlt_parallel/new_controller/include/new_controller/rolling_navigation.hpp|，第 99 行
\sourceentry{S50}{调度与保护}\par \path|../candidates/ldlt_parallel/new_controller/include/new_controller/rolling_navigation.hpp|，第 173 行
\sourceentry{S51}{ROS 输出}\par \path|../src/new_controller/src/navigation_input.cpp|，第 200 行
\sourceentry{S52}{FLARE 论文第 III 节、图 2–3 与式 (2)–(11)}\par \path|/home/starry/桌面/2027ICRA_FLARE.pdf|
\sourceentry{S53}{第 5 步实现和计时}\par \path|performance/step5_rolling_20261107/REPORT.md|
\sourceentry{S54}{R05 实验、计时和视频}\par \path|performance/r05_flight_20261107/REPORT.md|
\end{description}
\endgroup
\end{document}
